\begin{abstract}
\textbf{Objective:} Covariance-fitting fidelity,
predictive likelihood and source-participant
discrimination are distinct objectives for
EEG regularisation. We test when their
selection disagreement changes unknown-person
decoding, and whether a more exact source
search resolves the discrepancy.
\textbf{Methods:} We retain a source-fitted
covariance/tangent-LDA pipeline, participant-grouped
selection and source-only target information
boundaries. We correct a validation class-centred
covariance score to genuine source-trained-mean
Gaussian predictive likelihood, and decompose
AUC changes into positive--negative pair transitions.
Spectral interval bounds are used to search
continuous shrinkage ranges while reporting
conditional source-only optimality gaps.
\textbf{Results:} In a frozen EEGMMIDB
20-person block, source-discrimination selection
exceeded the historical class-centred score
by $0.0561$ AUC (Holm $p=0.00597$), although
five participants were harmed by more than
0.02 AUC. A three-channel BNCI2014\_004
worst-quarter source-risk criterion changed
eight of nine selections without a mean
AUC improvement. An anchored BNCI2014\_002
source search reduced an inconclusive
continuous-source optimality gap from
0.09477 to 0.01623 without improving
the original grid AUC. In a separate
protocol-first, unregistered Weibo2014
test (ten participants, 1,580 trials),
the continuous selector met its source
tolerance in all ten pools yet selected
the same shrinkage as the 27-point grid
in every pool; both obtained
0.73141 held-person macro AUC.
The gain relative to proper Gaussian
selection was only $0.00205$ AUC.
\textbf{Conclusion:} Closing a source
optimisation gap need not reduce transfer
risk. Correct validation estimands,
selection chronology and participant-level
negative results matter more than
shrinkage-grid precision alone.
\end{abstract}
